% -*- coding: utf-8 -*-
\chapter{Objetivos}

En este capítulo se definen las metas que marca la línea del desarrollo del proyecto. Se establece un objetivo general que engloba la finalidad última del asistente inteligente y un conjunto de objetivos específicos diseñados como metas tangibles, que garantizan el cumplimiento íntegro del proyecto, en base a las normativas vigentes.

\section{Objetivo general}
\begin{itemize}\item Diseñar, implementar y evaluar una plataforma de gestión de documentación para el sector agrícola basado en una arquitectura RAG multimodal, con el fin de analizar su impacto en la precisión, la relevancia, la trazabilidad y la reducción de alucinaciones en respuestas generadas sobre documentación agrícola y normativa. (Li, 2025)\end{itemize}
\section{Objetivos específicos}
\begin{itemize}\item Analizar el estado del arte sobre LLMs, RAG y alucinaciones en contexto regulados.\end{itemize}
\begin{itemize}\item Definir la arquitectura técnica del sistema y sus componentes principales.\end{itemize}
\begin{itemize}\item Construir una base de conocimiento vectorial a partir de documentación agrícola y normativa.\end{itemize}
\begin{itemize}\item Implementar un proceso de ingestación, fragmentación e indexación semántica del corpus documental.\end{itemize}
\begin{itemize}\item Evaluar el sistema mediante métricas de recuperación y generación.\end{itemize}
\begin{itemize}\item Comparar el rendimiento del LLM sin RAG frente al sistema propuesto con RAG.\end{itemize}


